
\subsubsection{2.1 Benchmark Saturation Studies}

The saturation of ML benchmarks has attracted growing attention. Bowman and Dahl [Bowman2021] argue that standard NLU benchmarks have largely saturated and call for principled retirement criteria --- but their analysis is qualitative; they identify the need for criteria without providing a quantitative threshold.

Liao et al. [Liao2022] conduct the largest systematic study of benchmark saturation, analyzing near-saturation trends across 3,765 benchmarks using time-based metrics and CoV. Their analysis confirms saturation is a systemic phenomenon --- but their metric is time-based (years to saturation), not paper\_count-based, and targets aggregate trends rather than discrete structural breaks.

Akhtar et al. [Akhtar2026] introduce the S\_index, a composite saturation metric for 60 LLM benchmarks, characterizing saturation as a continuous property. It does not apply change-point detection to PwC-internal CoV-vs-paper\_count data and does not produce a paper\_count threshold. Vasudevan et al. [Vasudevan2022] analyze ImageNet saturation in depth, documenting that top-1 accuracy above 90\% no longer discriminates between models --- providing independent evidence for the exploration-phase variance of pre-saturation benchmarks. Cao and Zhao [Cao2025] demonstrate that LLM benchmark scores are increasingly inflated by test-set pretraining, further motivating saturation detection.

\textbf{The gap:} No prior work applies change-point detection to PwC-internal CoV-vs-paper\_count data to produce a discrete, paper-count-based retirement threshold.

\subsubsection{2.2 Change-Point Detection Methods}

PELT (Pruned Exact Linear Time) [Killick2012] is an exact dynamic programming algorithm for detecting an unknown number of change-points in 1D signals with a penalized cost function. Wang, Lin, and Willett [Wang2019] establish the VPWBS algorithm's $O_p$(1/n) localization rate for regression change-points, providing theoretical grounding for PELT's performance at small sample sizes (N=115 in this work). The \texttt{ruptures} Python library [Truong2020] provides a well-maintained PELT implementation.

Existing benchmark analysis methods treat saturation as a continuous property or apply domain-specific heuristics. The approach in this paper instead treats the CoV-vs-paper\_count series as a signal in which a structural break should be detected and statistically validated.

\subsubsection{2.3 Benchmark Evaluation and Leaderboard Dynamics}

Papers With Code [PwC2019] maintains a comprehensive leaderboard database, making it a uniquely rich source for benchmark lifecycle analysis. Prior work using PwC data has examined research dynamics but not structural breaks in performance variability organized by publication count. CoV ($\sigma /\mu$ per benchmark) captures score diversity in a scale-invariant manner, suitable for pooling across heterogeneous metric scales [Liao2022].

This work is the first to apply PELT to PwC-internal residual CoV data, producing a quantitative paper\_count\* threshold grounded in change-point theory.

